% GENERATED FILE. Edit canonical lessons, then run npm run generate:books.
% canonical-source-hash: fnv1a64:a7f8dd4cf3207b7f
% canonical-lessons: TE-C338-kramanga, TE-C338-ranuranu, TE-C338-munupu, TE-C338-gatam, TE-C338-appatlo

\chapter{Little by Little}
\label{ch:te-little-by-little}
\hlchaptermodality{338}

\begin{chapteropening}
\textbf{By the end of this chapter:} \emph{I can say gradually, as time went by, previously, the past and in those days.}

Complete the last lesson of chapter 338: I can say gradually, as time went by, previously, the past and in those days.
\end{chapteropening}

\section[\texorpdfstring{\te{క్రమంగా} (kramaṅgā)}{క్రమంగా (kramaṅgā)}]{\te{క్రమంగా} (kramaṅgā) — gradually; step by step}

\label{lesson:TE-C338-kramanga}



\begin{quote}
Before the new one: say the Telugu for shackles, then the Telugu for the accused.
\end{quote}

\subsection*{You'll want to know: \te{క్రమంగా}}
\textbf{\te{క్రమంగా}} — \emph{kramaṅgā} — \textquotedblleft{}gradually; step by step\textquotedblright{}.

\begin{quote}
\te{క్రమంగా} \te{అతని} \te{ఆరోగ్యం} \te{బాగుపడింది}.
\end{quote}

\begin{quote}
\emph{kramaṅgā atani ārōgyaṁ bāgupaḍindi.}
\end{quote}

\begin{quote}
Gradually his health got better.
\end{quote}

\te{క్రమం}, order or sequence, with -\te{గా}: in order, step by step.

\begin{cousinweb}
\textbf{\te{క్రమంగా}} (\emph{kramaṅgā}) is built on Sanskrit \emph{krama}, \textquotedblleft{}a step, a series\textquotedblright{} (CDIAL 3577), with the Telugu adverb ending \emph{-gā}; CDIAL notes that Sanskrit \emph{kramēṇa} already meant \textquotedblleft{}by degrees\textquotedblright{}.
\end{cousinweb}

\begin{grammarlens}[title={What you've built}]
One more word for telling what happened.
\end{grammarlens}

\subsection*{Your turn}
\begin{itemize}
\raggedright
  \item \emph{Say it:} \emph{kramaṅgā}
  \item \emph{Say it:} \emph{kramaṅgā}, once more
  \item \emph{Say it:} \emph{kramaṅgā atani ārōgyaṁ bāgupaḍindi.} — Gradually his health got better.
  \item \emph{Recall:} say the Telugu for shackles, then the Telugu for the accused, then say \emph{kramaṅgā} again
\end{itemize}

\begin{tcolorbox}[breakable,colback=teal!4,colframe=teal!35!black,title={Before you move on}]
What does \te{క్రమంగా} mean? (\textquotedblleft{}Gradually; step by step\textquotedblright{}.) Say it once more.
\end{tcolorbox}
\section[\texorpdfstring{\te{రానురాను} (rānurānu)}{రానురాను (rānurānu)}]{\te{రానురాను} (rānurānu) — as time went by}

\label{lesson:TE-C338-ranuranu}



\begin{quote}
Before the new one: say the Telugu for the accused, then the Telugu for gradually.
\end{quote}

\subsection*{You'll want to know: \te{రానురాను}}
\textbf{\te{రానురాను}} — \emph{rānurānu} — \textquotedblleft{}as time went by\textquotedblright{}.

\begin{quote}
\te{రానురాను} \te{ఆ} \te{పల్లె} \te{పట్టణంగా} \te{మారింది}.
\end{quote}

\begin{quote}
\emph{rānurānu ā palle paṭṭaṇaṅgā mārindi.}
\end{quote}

\begin{quote}
As time went by, that village grew into a town.
\end{quote}

\te{రాను}, coming, said twice: as the days keep coming.

\begin{cousinweb}
\textbf{\te{రానురాను}} (\emph{rānurānu}) is native Telugu: the verb \emph{rā}, \textquotedblleft{}to come\textquotedblright{} (DEDR 5270), said twice.
\end{cousinweb}

\begin{grammarlens}[title={What you've built}]
One more word for telling what happened.
\end{grammarlens}

\subsection*{Your turn}
\begin{itemize}
\raggedright
  \item \emph{Say it:} \emph{rānurānu}
  \item \emph{Say it:} \emph{rānurānu}, once more
  \item \emph{Say it:} \emph{rānurānu ā palle paṭṭaṇaṅgā mārindi.} — As time went by, that village grew into a town.
  \item \emph{Recall:} say the Telugu for the accused, then the Telugu for gradually, then say \emph{rānurānu} again
\end{itemize}

\begin{tcolorbox}[breakable,colback=teal!4,colframe=teal!35!black,title={Before you move on}]
What does \te{రానురాను} mean? (\textquotedblleft{}As time went by\textquotedblright{}.) Say it once more.
\end{tcolorbox}
\section[\texorpdfstring{\te{మునుపు} (munupu)}{మునుపు (munupu)}]{\te{మునుపు} (munupu) — previously; before}

\label{lesson:TE-C338-munupu}



\begin{quote}
Before the new one: say the Telugu for gradually, then the Telugu for as time went by.
\end{quote}

\subsection*{You'll want to know: \te{మునుపు}}
\textbf{\te{మునుపు}} — \emph{munupu} — \textquotedblleft{}previously; before\textquotedblright{}.

\begin{quote}
\te{మునుపు} \te{ఇక్కడ} \te{ఒక} \te{చెరువు} \te{ఉండేది}.
\end{quote}

\begin{quote}
\emph{munupu ikkaḍa oka ceruvu uṇḍēdi.}
\end{quote}

\begin{quote}
Previously there was a pond here.
\end{quote}

Close to \te{ముందు}; \te{మునుపటి} means former: \te{మునుపటి} \te{రాజు}, the former king.

\begin{cousinweb}
\textbf{\te{మునుపు}} (\emph{munupu}) is native Dravidian: DEDR 5020 lists Telugu \emph{munupu}, \textquotedblleft{}former time; before\textquotedblright{}, with Tamil \emph{muṉpu}; \te{ముందు} comes from the same root.
\end{cousinweb}

\begin{grammarlens}[title={What you've built}]
One more word for telling what happened.
\end{grammarlens}

\subsection*{Your turn}
\begin{itemize}
\raggedright
  \item \emph{Say it:} \emph{munupu}
  \item \emph{Say it:} \emph{munupu}, once more
  \item \emph{Say it:} \emph{munupu ikkaḍa oka ceruvu uṇḍēdi.} — Previously there was a pond here.
  \item \emph{Recall:} say the Telugu for gradually, then the Telugu for as time went by, then say \emph{munupu} again
\end{itemize}

\begin{tcolorbox}[breakable,colback=teal!4,colframe=teal!35!black,title={Before you move on}]
What does \te{మునుపు} mean? (\textquotedblleft{}Previously; before\textquotedblright{}.) Say it once more.
\end{tcolorbox}
\section[\texorpdfstring{\te{గతం} (gataṁ)}{గతం (gataṁ)}]{\te{గతం} (gataṁ) — the past}

\label{lesson:TE-C338-gatam}



\begin{quote}
Before the new one: say the Telugu for as time went by, then the Telugu for previously.
\end{quote}

\subsection*{You'll want to know: \te{గతం}}
\textbf{\te{గతం}} — \emph{gataṁ} — \textquotedblleft{}the past\textquotedblright{}.

\begin{quote}
\te{ఆమె} \te{తరచుగా} \te{గతం} \te{గుర్తుచేసుకుంటుంది}.
\end{quote}

\begin{quote}
\emph{āme taracugā gataṁ gurtucēsukuṇṭundi.}
\end{quote}

\begin{quote}
She often recalls the past.
\end{quote}

You know \te{గత}, last or past: \te{గత} \te{వారం}. \te{గతం} is the noun.

\begin{cousinweb}
\textbf{\te{గతం}} (\emph{gataṁ}) is borrowed from Sanskrit \emph{gata}, \textquotedblleft{}gone, past\textquotedblright{} (CDIAL 4008); as a noun it is what has gone by.
\end{cousinweb}

\begin{grammarlens}[title={What you've built}]
One more word for telling what happened.
\end{grammarlens}

\subsection*{Your turn}
\begin{itemize}
\raggedright
  \item \emph{Say it:} \emph{gataṁ}
  \item \emph{Say it:} \emph{gataṁ}, once more
  \item \emph{Say it:} \emph{āme taracugā gataṁ gurtucēsukuṇṭundi.} — She often recalls the past.
  \item \emph{Recall:} say the Telugu for as time went by, then the Telugu for previously, then say \emph{gataṁ} again
\end{itemize}

\begin{tcolorbox}[breakable,colback=teal!4,colframe=teal!35!black,title={Before you move on}]
What does \te{గతం} mean? (\textquotedblleft{}The past\textquotedblright{}.) Say it once more.
\end{tcolorbox}
\section[\texorpdfstring{\te{అప్పట్లో} (appaṭlō)}{అప్పట్లో (appaṭlō)}]{\te{అప్పట్లో} (appaṭlō) — in those days}

\label{lesson:TE-C338-appatlo}



\begin{quote}
Before the new one: say the Telugu for previously, then the Telugu for the past.
\end{quote}

\subsection*{You'll want to know: \te{అప్పట్లో}}
\textbf{\te{అప్పట్లో}} — \emph{appaṭlō} — \textquotedblleft{}in those days\textquotedblright{}.

\begin{quote}
\te{అప్పట్లో} \te{ఇక్కడ} \te{బస్సులు} \te{లేవు}.
\end{quote}

\begin{quote}
\emph{appaṭlō ikkaḍa bassulu lēvu.}
\end{quote}

\begin{quote}
In those days there were no buses here.
\end{quote}

A spoken form: \te{అప్పట్లో} looks back at a whole period, \te{అప్పుడు} at one moment.

\begin{cousinweb}
\textbf{\te{అప్పట్లో}} (\emph{appaṭlō}) is native Telugu: \emph{appaṭi}, \textquotedblleft{}of that time\textquotedblright{}, from \emph{appuḍu}, \textquotedblleft{}then\textquotedblright{} (DEDR 4559), with the locative \emph{-lō}, \textquotedblleft{}in\textquotedblright{}.
\end{cousinweb}

\begin{grammarlens}[title={What you've built}]
One more word for telling what happened.
\end{grammarlens}

\subsection*{Your turn}
\begin{itemize}
\raggedright
  \item \emph{Say it:} \emph{appaṭlō}
  \item \emph{Say it:} \emph{appaṭlō}, once more
  \item \emph{Say it:} \emph{appaṭlō ikkaḍa bassulu lēvu.} — In those days there were no buses here.
  \item \emph{Recall:} say the Telugu for previously, then the Telugu for the past, then say \emph{appaṭlō} again
\end{itemize}

\begin{tcolorbox}[breakable,colback=teal!4,colframe=teal!35!black,title={Before you move on}]
What does \te{అప్పట్లో} mean? (\textquotedblleft{}In those days\textquotedblright{}.) Say it once more.
\end{tcolorbox}
